\documentclass{article}
\usepackage[T1]{fontenc}
\usepackage{amsmath}
\usepackage{amssymb}
\usepackage{longtable}
\usepackage[margin=2cm]{geometry}
\title{At rest}
\date{}
\begin{document}
\maketitle
\section*{Purpose and patterns}
Two bins at rest.

\section*{Entities, state variables and scales}
\subsection*{Bins}
\begin{longtable}{p{0.3\linewidth}p{0.3\linewidth}p{0.3\linewidth}}
Bin & Name & Unit \\
\hline
\endhead
A & A &  \\
B & B &  \\
\end{longtable}
\subsection*{Compartments}
\begin{longtable}{p{0.3\linewidth}p{0.3\linewidth}p{0.3\linewidth}}
Compartment & Bin & Categories \\
\hline
\endhead
A & A &  \\
B & B &  \\
\end{longtable}
\subsection*{Scales}
\begin{longtable}{p{0.45\linewidth}p{0.45\linewidth}}
Scale & Value \\
\hline
\endhead
Axis & one step is 1 step \\
Advance & continuous (StochasticEvents) \\
\end{longtable}
\section*{Process overview and scheduling}
\subsection*{Processes}
\begin{longtable}{p{0.225\linewidth}p{0.225\linewidth}p{0.225\linewidth}p{0.225\linewidth}}
Process & From & To & When \\
\hline
\endhead
ab & A, B & B, B & every step \\
in & none & A, B & every step \\
out & B & none & every step \\
dose & A & B & impulse \\
\end{longtable}
\subsection*{Scheduling}
Events happen one at a time, at random times: each event of a flow moves one unit from its sources to its targets, with a probability per unit of the axis (its propensity) of the rate times its source for a proportional rate, or the rate for an absolute one. The engine stops at each impulse, where the impulses act first, and at each position where a rate jumps.

\subsection*{Settings}
\begin{longtable}{p{0.45\linewidth}p{0.45\linewidth}}
Setting & Value \\
\hline
\endhead
method & tau\_leaping \\
tau & none \\
epsilon & 0.03 \\
tolerance & 1e-08 \\
rounding & none \\
\end{longtable}
\section*{Design concepts}
\begin{longtable}{p{0.45\linewidth}p{0.45\linewidth}}
Element & Description \\
\hline
\endhead
Basic principles & To be written by the modeller. \\
Emergence & To be written by the modeller. \\
Adaptation & To be written by the modeller. \\
Objectives & To be written by the modeller. \\
Learning & To be written by the modeller. \\
Prediction & To be written by the modeller. \\
Sensing & To be written by the modeller. \\
Interaction & A rate that reads neighbors\_\textless{}bin\textgreater{} or neighbors depends on every unit: every unit is a neighbor of every other. \\
Stochasticity & The engine draws the events at random. \\
Collectives & To be written by the modeller. \\
Observation & A run returns the value of every output at every position it reports; the submodels list the outputs. \\
\end{longtable}
\section*{Initialization}
\subsection*{Bin amounts}
\begin{longtable}{p{0.3\linewidth}p{0.3\linewidth}p{0.3\linewidth}}
Bin & Amount & Unit \\
\hline
\endhead
A & 5 &  \\
B & 5 &  \\
\end{longtable}
\subsection*{Initial state}
\begin{longtable}{p{0.3\linewidth}p{0.3\linewidth}p{0.3\linewidth}}
Compartment & Amount & Unit \\
\hline
\endhead
A & 5 &  \\
B & 5 &  \\
\end{longtable}
\section*{Input data}
None: the model reads no series of values.

\section*{Submodels}
\subsection*{Compact form}
\subsubsection*{Bin Transitions}
\noindent\texttt{ab (A, B -\textgreater{} B)}
\begin{align*}
A \to B &: 0 \cdot A \cdot B \\
B \to B &: 0 \cdot A \cdot B
\end{align*}
\noindent\texttt{in (none -\textgreater{} A, B)}
\begin{align*}
\varnothing \to A &: 0 \\
\varnothing \to B &: 0
\end{align*}
\noindent\texttt{out (B -\textgreater{} none)}
\begin{align*}
B \to \varnothing &: 0 \cdot B
\end{align*}
\noindent\texttt{dose (A -\textgreater{} B, impulse at 1)}
\begin{align*}
A \to B &: 0 \cdot A
\end{align*}
Total System: 2 coupled equations (2 bins)

\subsection*{Events}
\begin{longtable}{p{0.3\linewidth}p{0.3\linewidth}p{0.3\linewidth}}
Transition & Change & Propensity \\
\hline
\endhead
ab & A -1, B +1 & 0 * A * B \\
in & A +1, B +1 & 0 \\
out & B -1 & 0 * B \\
\end{longtable}
Each event moves one unit. A flow has no events while one of its sources holds fewer units than an event takes.

\subsection*{Impulses}
\noindent\texttt{dose, impulse at 1}
\begin{align*}
A \to B &: 0 \cdot A
\end{align*}
\subsection*{Neighbors}
Every unit is a neighbor of every other (Complete): neighbors\_\textless{}bin\textgreater{} is the bin total and neighbors is total.

\subsection*{Parameters}
The model has no parameters.

\subsection*{Outputs}
\begin{longtable}{p{0.45\linewidth}p{0.45\linewidth}}
Source & Outputs \\
\hline
\endhead
A & A \\
B & B \\
\end{longtable}
\end{document}
